\section{Books, Revision, and the Long Accident of Survival}

\subsection{A career distributed across arguments}

Tertullian left more than thirty works generally received as authentic, but an
exact number changes with decisions about fragments, titles, composite texts,
and attribution. They range from brief addresses to five-book controversy. No
ancient collected edition preserves their authorial order, and the works rarely
name a year. Chronology is reconstructed from emperors and governors, eclipses,
allusions among books, doctrinal vocabulary, positions on discipline, and
manuscript grouping. Circular dating is a constant danger: a work should not be
dated ``Montanist'' only because it is rigorous and then used to prove that
rigor belongs to a dated Montanist phase.

The corpus includes several different acts of speech:

\begingroup
\small
\setlength{\tabcolsep}{4pt}
\begin{longtable}{>{\raggedright\arraybackslash}p{0.20\linewidth}
                  >{\raggedright\arraybackslash}p{0.28\linewidth}
                  >{\raggedright\arraybackslash}p{0.44\linewidth}}
\toprule
\textbf{Work family} & \textbf{Representative works} &
\textbf{Biographical and textual control}\\
\midrule
\endhead
Public apologetic & \emph{To the Nations}, \emph{Apologeticum}, \emph{To
Scapula} & Addresses imagined pagan publics and named authority. Forensic
claims, reported pagan speech, demographic boasting, and threats of mass
confession are rhetoric as well as evidence.\\
Martyrdom and civic separation & \emph{To the Martyrs}, \emph{Scorpiace},
\emph{On Flight in Persecution}, \emph{On Idolatry}, \emph{On the Crown},
\emph{On Spectacles} & Different occasions and changing judgments about flight,
service, trades, and confession must be compared rather than harmonized into a
single rule.\\
Worship and formation & \emph{On Baptism}, \emph{On Prayer}, \emph{On
Repentance}, \emph{On Patience} & Treatises persuade and correct actual or
imagined audiences; practice described in argument is not automatically a
universal ritual rubric.\\
Household and sexual discipline & \emph{To His Wife}, \emph{On the Apparel of
Women}, \emph{On the Veiling of Virgins}, \emph{Exhortation to Chastity},
\emph{On Monogamy} & The sequence exposes both continuity and hardening.
Addresses to women and a wife are evidence of relation, not permission to
invent the addressees' responses.\\
Rule and anti-heretical polemic & \emph{Prescription against Heretics},
\emph{Against Hermogenes}, \emph{Against the Valentinians}, \emph{Against
Marcion}, \emph{Against Praxeas}, \emph{On the Flesh of Christ} & Opponents'
teachings are mediated by adversarial selection. \emph{Against Marcion} has an
explicit multi-edition history; doctrinal works also respond to particular
arguments rather than supplying timeless summaries.\\
Anthropology and resurrection & \emph{On the Soul}, \emph{On the Resurrection
of the Flesh} & Scriptural exegesis is joined to Stoic, medical, philosophical,
and visionary material. Later labels must preserve Tertullian's distinctive
meanings.\\
Prophetic discipline & \emph{On Fasting}, \emph{On Modesty}, \emph{On
Monogamy}, and New Prophecy passages across other works & This is not a sealed
late shelf: allegiance and discipline cross genres, and individual dates remain
disputed.\\
\bottomrule
\end{longtable}
\endgroup

The author is sometimes an unreliable witness against himself in a productive
sense. \emph{On Patience} 1 says he is wholly unfit to teach the virtue he does
not possess and can only sigh after it like a sick person describing health.
The confession is a literary humility topos and recognizable self-criticism.
It establishes that Tertullian wanted readers to see the gap between teacher
and achievement; it does not license a diagnosis of temperament from every
insult.

\subsection{Latin originality and missing Greek}

The surviving corpus is overwhelmingly Latin, but Tertullian moved in a
bilingual Christian intellectual world. \emph{On Baptism} 15 refers to a
fuller treatment already composed in Greek. \emph{On the Crown} 6.3 says his
related treatment of spectacles was in Greek, and \emph{On the Veiling of
Virgins} 1.1 says he had already handled that question in Greek. Those
authorial Greek treatments are lost. They must be distinguished from the
anonymous ancient Greek translation of the Latin \emph{Apologeticum} used by
Eusebius, whose translator and complete text do not survive. Other lost titles
are known from Tertullian's cross-references or later lists, including works on
ecstasy closely tied to the New Prophecy.

Counts of fifteen, eighteen, or another exact total of lost works depend on
whether alternate titles, Greek versions, projected works, and fragments are
separate. The responsible conclusion is qualitative: substantial parts of his
Greek production, prophetic defense, and intra-Christian controversy are
missing. The surviving Latin career is therefore not the whole author.

His Latin is itself difficult evidence. Tertullian compresses clauses, coins or
retools terms, shifts registers from street taunt to technical argument, and
exploits sounds and legal ambiguity. A smooth translation can conceal the
argument; a wooden one can manufacture obscurity. Claims about the first use of
a theological word require the Latin, earlier lexical record, textual variant,
and sense in context.

\subsection{Four corpus routes and one independent apology}

No surviving manuscript is close to Tertullian's lifetime. The earliest located
item is a late-eighth-century fragment of \emph{Against the Jews} in Paris,
Biblioth\`eque nationale de France, lat. 13047. The large manuscript witnesses
begin still later. Most surviving copies are medieval or Renaissance, and many
works depend on very few witnesses.

Modern editors describe several principal collection histories rather than one
stemma for all works:

\begin{itemize}
\item The \emph{Apologeticum} circulated independently and survives in more
than three dozen manuscripts, with textual groupings of its own.
\item The \emph{Corpus Cluniacense} transmitted a large collection of roughly
twenty-one or twenty-two works through branches often called alpha and beta.
The counts and relationships vary with reconstructed archetypes and losses.
\item The \emph{Corpus Agobardinum} is represented above all by the
ninth-century Codex Agobardinus. It preserves thirteen works from a larger
collection and is uniquely important for several readings.
\item The \emph{Corpus Trecense} survives in a sole extant representative, the
twelfth-century Troyes manuscript 523, which carries five works of Tertullian.
That family-level fact does not by itself make each work a codex unicus; every
treatise's other transmission must still be checked separately.
\item The \emph{Corpus Corbeiense} carried a group of major anti-heretical
works. Its manuscript witnesses were lost, apart from a fragment and old
collations; its reconstructed contents and readings depend partly on those
indirect controls and on sixteenth-century editions made from codices no longer
available.
\item An Ottobonian excerpt collection and other partial witnesses preserve
additional pieces. Excerpt is not the same object as a complete treatise.
\end{itemize}

These routes explain why one general statement about ``the manuscripts of
Tertullian'' is unsafe. A richly witnessed \emph{Apologeticum}, a sole-copy
work, and a text reconstructed partly from early print do not carry the same
security at every word. The printed editions of Beatus Rhenanus, Mesnart,
Pamelius, and others can be textual witnesses where their exemplars vanished,
even though editorial conjecture and printing error must also be controlled.

The 1954 \emph{Corpus Christianorum, Series Latina} volumes 1--2 remain a
standard compact critical corpus, but they contain work-specific editions by
different editors rather than one uniform new collation. The \emph{Sources
Chr\'etiennes} series and other modern editions have since re-edited individual
works with fuller introductions, manuscript arguments, translations, or
commentary. Migne and searchable transcriptions are useful discovery tools; no
critical claim should rest on their convenience alone.

\subsection{Authentic, appended, attributed}

The short \emph{Against All Heresies}, often printed after \emph{Prescription
against Heretics}, is now treated as Pseudo-Tertullian. Its placement in the
transmitted corpus records reception, not authorship. The verse \emph{Poem
against Marcion} and several other works transmitted under his name are likewise
spurious. A maxim or theological position in one of them cannot be moved into
Tertullian's biography without an attribution argument.

Attribution also fails in subtler ways. Tertullian's proposed editorship of the
\emph{Passion of Perpetua and Felicity} rests on affinities of language,
martyrdom, prophecy, and the text's editorial frame. Similarity may reflect a
shared Carthaginian Christian world. It does not identify an editor. The work
belongs in his context and reception, not securely in his bibliography.

The uneven survival changes interpretation. A lost Greek treatise makes the
Latin author seem more monolingual; lost books on ecstasy make his prophetic
case look like scattered discipline; spurious anti-heretical catalogs can make
his taxonomy seem larger; early print can freeze one reconstructed reading.
Transmission is therefore not a technical epilogue to the life. It is one of
the forces that made the Tertullian later centuries could praise or condemn.
